\section{Introduction}
\label{sec:intro}

A driver should be able to say \vi{bật điều hoà 22 độ}{turn the AC on at 22 degrees} or simply \vi{trong xe bí quá}{it is
stuffy in here} and have the car respond, without taking their hands off the wheel and without depending on a mobile
connection. Large language models (LLMs) can map such utterances to structured function calls, but the usual recipe---placing
the schemas of every available function in the prompt of a large model---is slow and memory-hungry on the kind of hardware
found in a car, and requires a network connection when the model is hosted remotely.

Octopus~v2 \citep{chen2024octopus} proposed a different recipe for on-device agents: each function becomes a new
\emph{functional token}, a small model is fine-tuned to emit \tok{nexa\_i}\code{(arguments)}, and the function descriptions
are only needed at training time. Function selection then reduces to predicting one token, and the inference prompt contains
little more than the user's sentence. Follow-up work applied small models to in-vehicle function calling
\citep{sowti2025optimizing,jafarzadeh2026fixedkeys} and showed a clear trade-off: functional tokens are many times faster than
schema-in-prompt, but tend to call a function when they should decline \citep{jafarzadeh2026fixedkeys}.

This report describes \textbf{Nexa Drive}, a graduation project that builds such an assistant for \emph{Vietnamese} and takes
it all the way to a working product: speech recognition, a wake word, runtime guardrails, speech synthesis and an interactive
3D car. The product is the goal; the research layer exists to justify each technical choice with measurements. The setting
raises four problems that published function-calling work rarely addresses together:
\begin{enumerate}
  \item \textbf{No Vietnamese in-car data.} We generate all training data with an open 72B teacher run locally on rented
  GPUs (no paid API), and evaluate on separately generated, human-reviewed sets that are frozen by hash before any model is
  scored.
  \item \textbf{Drivers rarely speak in commands.} Many requests are complaints (\vi{nóng quá}{it is so hot}) whose right
  answer is neither a call nor a refusal but a proposal. We add a third action, \emph{clarify}, marked by a
  \tok{nexa\_ask} token, and measure it on dedicated sets.
  \item \textbf{A wrong action moves a physical car.} We keep the model responsible for understanding and put safety in
  code: a confidence threshold on the decision token, a parameter validator, state-dependent rules (no opening the trunk
  while driving) and confirmation for sensitive actions or wake-word commands.
  \item \textbf{Speech is the main interface.} Errors of the speech recognizer, Vietnamese number words, background noise
  and false wake-ups all reach the language model, so we evaluate the whole chain with real recordings.
\end{enumerate}

\paragraph{Contributions.}
\begin{itemize}
  \item A complete on-device Vietnamese in-car assistant: a 1.5B functional-token model (1.04\,GB after 4-bit
  quantization, median 0.52\,s per decision on a laptop CPU), speech recognition, speech synthesis, a custom wake word,
  runtime guardrails and a 3D interface, all running offline (\S\ref{sec:system}, \S\ref{sec:speech}).
  \item A distillation pipeline with code-first validation, a two-family LLM jury and human audits, together with a record of
  what did and did not work over six data rounds---most notably that hard examples must come in \emph{contrastive pairs}
  across labels, or the model learns a shortcut from sentence form (\S\ref{sec:data}, \S\ref{sec:rounds}).
  \item An evaluation protocol for small test sets: frozen and contamination-checked sets, two training seeds per
  configuration, exact McNemar tests and acceptance criteria written down before each run. Under this protocol several
  apparent improvements turned out to be seed noise, and we report them as such (\S\ref{sec:protocol}).
  \item Measured answers to practical questions: how a 1.5B student compares with its 72B teacher (\S\ref{sec:main}), LoRA
  versus full fine-tuning (\S\ref{sec:lora}), what a confidence threshold buys (\S\ref{sec:threshold}), what 4-bit
  quantization costs (\S\ref{sec:quant}), whether GRPO helps on distilled data (\S\ref{sec:grpo}), and how much speech
  recognition and noise degrade the chain (\S\ref{sec:speech}). Negative results are reported alongside positive ones.
\end{itemize}

The vehicle is simulated: the assistant drives a state machine and a 3D model rather than a real CAN bus. All evaluation sets
are synthetic (generated by an LLM from a written specification and reviewed by a human), and all voice recordings come from
one speaker. We return to these limitations at the end of the report.
